\addcontentsline{toc}{chapter}{Abstract}\vspace{-1cm}
%Border
\begin{tikzpicture}[remember picture, overlay]
  \draw[line width = 4pt] ($(current page.north west) + (0.75in,-0.75in)$) rectangle ($(current page.south east) + (-0.75in,0.75in)$);
\end{tikzpicture}

Transmission lines are critical components of modern power systems, responsible for delivering electricity from generation units to consumers. However, they are highly susceptible to faults caused by environmental disturbances, equipment failures, and load variations. These faults can lead to system instability, equipment damage, and large-scale power outages if not detected and handled promptly. Traditional fault detection methods, primarily based on relay systems, often lack adaptability and fail to provide deeper insights into system-wide impacts. With the increasing complexity of smart grids, there is a growing need for intelligent, automated, and real-time monitoring solutions.

This report presents an advanced intelligent fault detection and predictive governance system that integrates Embedded Machine Learning, Decision Intelligence, and Generative AI. The system utilizes real-time voltage and current sensor data to detect and classify faults using a lightweight Decision Tree model deployed on embedded hardware. Beyond detection, a Decision Intelligence Engine analyzes cascading effects and evaluates system-wide impact, enabling proactive risk assessment. A Llama 3-based Generative AI layer is incorporated to perform root cause analysis, generate natural language explanations, and suggest mitigation strategies. Additionally, a Smart Governance Layer automates alerts, escalation workflows, and policy-driven responses, ensuring efficient and structured fault management.

The proposed system demonstrates high fault classification accuracy with low latency suitable for real-time applications. The integration of decision intelligence and AI-driven interpretation significantly enhances system usability and operational efficiency compared to traditional approaches. Furthermore, the system introduces a predictive governance paradigm by mapping fault events to potential operational and economic impacts, enabling better decision-making. The overall architecture provides a scalable and practical solution for intelligent power system monitoring, with potential applications in smart grid management and industrial energy systems.

\pagebreak